\section{Introduction}
\label{sec:intro}

A learner that takes supervision from other learners cannot treat all of its
sources alike. Some are better than others, some are adversarial, and the
standard response is to estimate a per-source reliability and weight incoming
gradients by it. This is one of the most widely deployed ideas in machine
learning: it is the Dawid--Skene construction in crowdsourcing, annotator
modelling in preference learning, client reweighting in federated optimisation,
and peer weighting in decentralised training. In each case reliability estimation
is introduced defensively --- it is the component that makes learning
\emph{robust} to bad sources.

This paper reports a cost of that component that we have not seen stated. A
per-source reliability estimate is a variable indexed by the source. If anything
about how a learner encodes its inputs depends on \emph{which} source sent them,
the reliability estimate is the natural place for that dependence to accumulate.
When the source-dependent part of the encoding happens to correlate with a group
label, reliability estimation converts a local encoding defect into a global,
self-reinforcing partition of the population --- and does so with a sharp
threshold rather than gracefully.

We make this precise in a population of $N$ agents that exchange opinions about a
fixed agenda of $P$ issues. Each agent carries a classifier (the \emph{opinion}
sector) and, for every other agent, a scalar reliability estimate (the
\emph{trust} sector). Both are updated by the same rule, a Gaussian variational
approximation to the Bayesian posterior that is optimal for a single
teacher--student pair. The perturbation we inject is a representation error in
the sense familiar from shortcut learning: a discriminating agent appends to each
message a feature that is uninformative about the issue and determined entirely by
the sender's group label. It shifts one field by a constant $d$. Nothing in the
model expresses a preference for either group, no training data is biased, and
the label is uncorrelated with every issue.

\paragraph{Contributions.}
\begin{itemize}
  \item We identify \textbf{learned per-source reliability weighting as an
    amplifier of spurious features}. A perturbation that is harmless when keyed to
    a random per-source bit produces a population-wide split when keyed to a group
    label of the same magnitude (\S\ref{sec:ablation}, condition A0).
  \item We map the \textbf{phase diagram} in the strength $d$ of the spurious
    feature and the fraction $f_d$ of agents carrying it, extract the boundary
    $d_c(f_d)$, and separate \textbf{two distinct discriminatory regimes}: one in
    which the split still tracks learned representations, and one in which trust
    tracks the group label almost independently of them
    (\S\ref{sec:phase}).
  \item We \textbf{isolate the mechanism by ablation}. Freezing the trust sector,
    or removing the sign reversal by which distrusted sources are anti-learned,
    removes the transition; Hebbian and Perceptron updates under the identical
    perturbation do not exhibit it (\S\ref{sec:ablation}).
  \item We show that \textbf{task diversity $\alpha = P/K$ determines the order in
    which structure crystallises} --- trust before representation, or the reverse
    --- which is a statement about what a population of learners aligns first, and
    is independent of whether any spurious feature is present at all
    (\S\ref{sec:agenda}).
\end{itemize}

\paragraph{The frame.} The update rule we study is optimal in a well-defined
sense, for a pair. Deployed in a population it is pathological, and the
pathology is not a bug in the derivation --- it follows from the two features that
make the rule good at its stated job. We read this as objective misspecification
by scale: an objective correct at the scale it was derived for, wrong at the scale
it is used at. We state this once here and once in \S\ref{sec:limitations}, and
otherwise let the ablations carry the argument.

\todo{one short paragraph on why this is not merely a restatement of model
collapse --- there the degradation is in the data distribution, here the data
distribution is fixed and the pathology is in the weighting.}
